\section{Property Implication Analysis}
\label{sec:property}

% Section 5 of the paper.
% Phase 6: properties as first-class objects
% Universal detection via double negation
% implies_observed / equivalent_observed
% Full table: 10 axiom classes x 3 property pairs

\subsection{Properties as First-Class Objects}

% define_property API, has_property_1, is_property
% Detection (forward): formula -> has_property
% No backward application (cascading prevention)

\subsection{Universal Property Detection}

% Existential detection insufficient for algebraic properties
% Double negation: not_left_id(e) <- element(e), element(x), NOT op(e,x,x)
%                  left_id(e) <- element(e), NOT not_left_id(e)

\subsection{Implication Discovery}

% derive_property_implications: ext(P) subset ext(Q) -> implies_observed(P,Q)
% implies_observed vs implies (inductive vs deductive)
% Modus ponens only on deductive implies

\subsection{Full Property Table}

% Table 1: 10 axiom classes x {left<->right, left<->idem, right<->idem}

\begin{table}[h]
\centering
\caption{Property equivalences across axiom classes on 3-element carriers.
$\leftrightarrow$ indicates observed equivalence (bidirectional implication);
$\times$ indicates no equivalence. All properties detected via universal
(double-negation) mechanism.}
\label{tab:property-table}
\begin{tabular}{lccc}
\toprule
Axiom class & left $\leftrightarrow$ right & left $\leftrightarrow$ idem & right $\leftrightarrow$ idem \\
\midrule
\texttt{assoc}             & $\times$ & $\times$ & $\times$ \\
\texttt{assoc+comm}        & $\times$ & $\times$ & $\times$ \\
\texttt{assoc+comm+id}     & $\leftrightarrow$ & $\times$ & $\times$ \\
\texttt{assoc+comm+id+inv} & $\leftrightarrow$ & $\leftrightarrow$ & $\leftrightarrow$ \\
\texttt{assoc+id}          & $\leftrightarrow$ & $\times$ & $\times$ \\
\texttt{comm}              & $\times$ & $\times$ & $\times$ \\
\texttt{comm+id}           & $\leftrightarrow$ & $\leftrightarrow$ & $\leftrightarrow$ \\
\texttt{comm+id+inv}       & $\leftrightarrow$ & $\leftrightarrow$ & $\leftrightarrow$ \\
\texttt{id}                & $\leftrightarrow$ & $\times$ & $\times$ \\
\texttt{id+inv}            & $\leftrightarrow$ & $\leftrightarrow$ & $\leftrightarrow$ \\
\bottomrule
\end{tabular}
\end{table}

% Key findings:
% 1. id present <-> left=right (necessary and sufficient)
% 2. inv or comm+id -> identity = unique idempotent
% 3. assoc alone does not affect identity equivalence
